\section{Introduction}
\label{sec:intro}

Language-driven grasp detection requires predicting a five-parameter grasp
rectangle
\begin{equation}
  G=(x, y, w, h, \theta)
  \label{eq:grasp}
\end{equation}
for the object or part selected by a free-form instruction.  We start from the
diffusion-based LGDM formulation.  In its released code, the vision-language
fusion output is computed but not injected into the grasp decoder, and the
diffusion loss is not optimized.  We define a clean baseline that restores the
diffusion objective and introduce LSAR, a compact residual module that makes
the language condition spatially explicit.  The main contribution is the
LSAR conditioning mechanism; the repaired baseline, repeated evaluations, and
ablation are provided to make the effect reproducible and interpretable.

\section{Related Work}
Grasp detection methods based on dense quality, angle, and width maps
\cite{morrison2018closing,kumra2020grconvnet} provide the representation used
by many language-conditioned systems.  Language-driven Grasp Detection (LGD)
\cite{vuong2024lgd} introduces Grasp-Anything++
\cite{vuong2023graspanything} and a diffusion-based conditional model.
Vision-language encoders such as ALBEF \cite{li2021albef} align semantic
concepts but do not directly encode graspable regions.  Our module adapts the
language-conditioned spatial branch without changing the pretrained encoders
or the generative head.
